\usepackage{xfakebold,tcolorbox,amsmath,amssymb,graphicx,ascmac,fancybox,framed,tikz,comment,calc,amsthm,bm,enumitem}
\tcbuselibrary{breakable, skins, theorems}
\usetikzlibrary{positioning, intersections, calc, arrows.meta,math}
\usepackage[top=11truemm,left=20truemm,right=15truemm,bottom=23truemm]{geometry}

% \white :計算用余白の挟み込み
\newcommand{\white}{\newgeometry{top=20truemm}\leavevmode
\begin{center}
    \noindent(\,計　算　用　余　白\,)
\end{center}}
% \exc[サイズ]{数値} :a\exc{b}でa^b, 指数に指数を載せる時に使用　数式環境内ではサイズ0.95
\newcommand{\exc}[2][0.75]{^{\textstyle\scalebox{#1}[#1]{\raise0.25ex\hbox{$#2$}}}}
% \exfrac[サイズ]{分母}{分子} :指数部分に分数を持ってくるときに使う．積分範囲などに使う場合サイズ0.75~0.8を推奨
\newcommand{\exfrac}[3][0.90]{\tfrac{\lower0.1ex\hbox{\scalebox{#1}[#1]{$#2$}}}{\raise0.3ex\hbox{\scalebox{#1}[#1]{$#3$\rule{0pt}{7.5pt}}}}}
% \dint[文字サイズ]{下付}{上付} :インテグラルのdisplay省略、積分範囲の表示サイズを変更、
\newcommand{\dint}[3][0.75]{\displaystyle\int_{\scalebox{#1}[#1]{$#2$}}^{\scalebox{#1}[#1]{$#3$}}}
% \dsqrt[n乗根]{中身}　:√を重ねる時用
\newcommand{\dsqrt}[2][]{\displaystyle\sqrt[\raise0ex\hbox{{\scalebox{0.625}[0.625]{$#1$}}}]{#2}}
% \dlim{下に付けるヤツ} :文字サイズ変更に適用させたlim
\newcommand{\dlim}[1]{\raise0.3ex\hbox{\scalebox{1.04}{${\normalsize \displaystyle\lim_{#1}}$}}\,}
% \cmb{n}{r}  \prm{n}{r}　: nCr nPr
\def\cmb#1#2{_{#1}\mathrm{C}_{#2}}\def\prm#1#2{_{#1}\mathrm{P}_{#2}}
% \w :(n)の後の空白用
\def\w{\;\;\;}\def\spw{\;\;}
% \hs \vs \vhs  :未使用
\def\hs{\hspace{1.38cm}}\def\vs{\vspace{-2mm}\\[1mm]}\def\vhs{\vs\hs}
\def\mrm{\mathrm}
\def\ve{\overrightarrow}
% 表紙のページ番号を省略
\setcounter{page}{0}